% Exercise splitting, nested math/list content, and proof end marks.
\documentclass[oneside,openany]{bookbangla}
\setbanglamainfont{Kalpurush}
\title{বাক্সের পরীক্ষা}
\newcount\repeatcount
\newcommand\longpassage{%
  \repeatcount=0
  \loop
    একটি দীর্ঘ অনুচ্ছেদ। বাংলা ও English লেখা পাশাপাশি থাকবে।
    পরের পাতায় গেলেও বাক্সের লেখা ও সীমানা ঠিক থাকতে হবে।\par
    \advance\repeatcount by1
  \ifnum\repeatcount<24\repeat}
\begin{document}
\mainmatter
\chapter{দীর্ঘ লেখা}
\begin{example}[label=example:long]\longpassage\end{example}
\begin{reminder}[বিশেষ কথা]\longpassage\end{reminder}
\begin{diybox}\longpassage\end{diybox}
\begin{exercise}[hint={এখানে $x^2+y^2$ এবং \textbf{একটি সূত্র} আছে।}]
\label{exercise:long}
\longpassage
\begin{enumerate}\item প্রথম ধাপ।\item দ্বিতীয় ধাপ।\end{enumerate}
\end{exercise}
\begin{solution}[অন্য সমাধান]
শেষে একটি সমীকরণ।
\begin{equation*}1+1=2.\qedhere\end{equation*}
\end{solution}
\begin{proof}একটি প্রমাণ।\end{proof}
\autoref{exercise:long}; \autoref{example:long}.
\end{document}
